\chapter{Modello di dominio}
\label{ch:dominio}

\section{Topic}
\texttt{Topic(id, title, isOptional, status, subtopics)}.
\texttt{status} $\in$ \{locked, inProgress, completed\}. I sottotopic ereditano
la visibilità click-to-expand (US-01). I topic opzionali sono evidenziati ma
non bloccano l'avanzamento.

\section{Quiz}
\texttt{Quiz(topicId, questions[5], threshold=0.8)}.
Soglia di superamento: $required = \lceil 0.8 \times n \rceil$; con $n=5$,
$required = 4$. Ogni risposta dà feedback immediato (corretto/errato);
il superamento sblocca il topic successivo e assegna la scelta ricompensa
(US-02/US-03).

\section{Reward (carta)}
\texttt{Reward(id, name, damage, effect, rarity, rarityColor)}.
\texttt{rarity} $\in$ \{common, rare, epic, legendary\}; \texttt{rarityColor}
è una stringa esadecimale (es.\ \texttt{\#9E9E9E}). Effetti noti: danno
diretto, cura, veleno. Il giocatore ne sceglie 1 per topic; l'inventory
è il deck (max 6 carte attive + 2 slot bloccati da sbloccare).

\section{BossFight}
\texttt{BossFight(boss, playerHp, thresholds=\{25,50,75\}\%)}.
Parametri consolidati dal prototipo:
\begin{itemize}[nosep]
\item danno carta base 15 + \texttt{damage} della carta;
\item attacco normale boss 10, speciale 20, cura boss $-15$ (al boss),
veleno 5/turno;
\item danno da quiz: $\mathit{round}(\mathit{percentuale}/10)$;
\item alle soglie 25/50/75\% il boss cambia pattern (attacco speciale o cura).
\end{itemize}
Esiste un \texttt{TODO} noto (\texttt{victory} in \texttt{boss\_fight\_active},
riga $\sim$419) sulla gestione della vittoria.

\section{PlayerProgress}
\texttt{PlayerProgress(xp, level=1, completedTopics, deck, bossStates)}.
Ogni level-up assegna +100\,XP ma il level è attualmente fisso a 1
(manca la curva di crescita). Nota: \texttt{fromJson} perde il campo
\texttt{adaptiveQuizzes} (debito noto, Cap.~\ref{ch:debiti}).
